\chapter*{Preface}
\addcontentsline{toc}{part}{Preface}



\section*{Structure of dissertation}

The doctoral dissertation is devided in four parts:
\begin{itemize}
	\item In part \hyperref[part:I]{\Romanbar{I}} the reader is introduced to the topic of the dissertation.
	The motivation of research carried out in connection with the doctoral dissertation is presented.
	The reader is acquainted with the current problems of integration optoelectronics with conventional electronics as well as the requirements that optoelectronic components materials must meet.
	The conception of band structure engineering is introduced along with short explanation of main methods and their practical implementations.
	Finally, the theoretical tools essential for research are introduced.
	\item Part \hyperref[part:II]{\Romanbar{II}} concerns the methodology of conducted research and a brief description of tools used to obtain results.
	The \acrshort{DFT} based \acrshort{abinitio} calculations and the theoretical background are introduced providing basic understanding of the physics behind the codes along with advantages and disadvantages of implemented solutions.
	The band structure engineering methods are shown in detail.
	Under each of them, the whole algorithms are characterized and every step of the sequence is thoroughly explained and the usage justified.
	\item In part \hyperref[part:III]{\Romanbar{III}} the outcome of the research is presented along with the analysis. The results were also published in scientific journals:
	\begin{itemize}
		\item [\cite{janik.2020}] Norbert Janik, Pawel Scharoch and Robert Kudrawiec,  Towards band gap engineering via biaxial and axial strain in group IV crystals, \href{https://doi.org/10.1016/j.commatsci.2020.109729}{\textsw{Computational Materials Science}}, \href{https://www.sciencedirect.com/science/article/pii/S0927025620302202}, 2020;
		\item [\cite{janik.2021}] Paweł Scharoch, Norbert Janik, Michał Wiśniewski, Herbert S. Mączko, Marta Gładysiewicz, Maciej P. Polak and Robert Kudrawiec,  Electronic band structure of semiconductor alloys: From ab initio to $k\cdot p$ via computational alchemy, on example of $\textrm{Ge}_{1-x}\textrm{Sn}_x$ alloy, \href{https://doi.org/10.1016/j.commatsci.2020.110052}{\textsw{Computational Materials Science}}, \href{https://www.sciencedirect.com/science/article/pii/S0927025620305437}, 2021;
	\end{itemize}
	\item Part \hyperref[part:IV]{\Romanbar{IV}} consists of a specific supplement to the doctoral dissertation below.
	At the end of this part there are lists of: acronyms, symbols, figures and a bibliography. 
\end{itemize}

\ornament


\section*{Aims of work}

The objectives of this doctoral dissertation are listed below. 
\begin{itemize}

    \item Conduct wide, full scale research of band structure engineering methods of group IV materials and the changes that those approaches provide. Also giving a complete picture of possible properties that can emerge for modifications of group IV materials.
    
    \item Find candidates for materials applicable in optoelectronics that at the same time will allow for easy integration with conventional electronics.
    
    \item Development of a method to efficiently analyze results of calculations of each type of band structure engineering. Algorithms will be capable to implement for almost any kind of materials group, providing a tool to investigate possibly achievable different electronic properties.
\end{itemize}

\ornament